% !TeX root = ../main.tex
\chapter{The No-Trade Theorem}
\chapterauthor{Aden Chen}

% !TeX root = ../main.tex
\section{Farkas' Lemma}
This section focuses on Farkas' Lemma, an important result concerning systems of linear equations.
For notational clarity, we will follow the tradition of generally not explicitly indicating transposes.

There are many mathematically equivalent ways to state Farkas' Lemma.
We start with two versions with particularity nice geometric and algebraic interpretations.

\begin{lemma}[Farkas' Lemma, Geometric version]\label{lem:farkas-geom}
  Fix \(A \in \bF^{m \times n}\) and \(b \in \bF^m\).
  Exactly one of the following is true:
  \begin{itemize}
    \item There exists \(x \in \bF^n\), \(x \geq 0\)\footnote{\(\geq 0\) is interpreted as component wise.} such that \(Ax = b\).
    \item There exists \(y \in \bF^m\) such that \(yA \geq 0\) and \(y b < 0\).
  \end{itemize}
\end{lemma}

This says that either \(b\) lies in the positive cone generated by columns of \(A\) or there is a hyperplane separating this cone from \(b\).
\mnote{
  Thus one can prove the result above in the case of \(\bR\) using some hyperplane separation theorem for convex sets, which typically relies on topological properties of the field.
  However, Farkas' Lemma remains true for arbitrary ordered fields.
  For those cases, one needs to use an algebraic argument.
}

\begin{lemma}[Farkas' Lemma, Algebraic version]\label{lem:farkas-alg}
  Fix \(A \in \bF^{m \times n}\) and \(b \in \bF^m\).
  Exactly one of the following is true:
  \begin{itemize}
    \item There exists \(x \in \bF^n\) such that \(Ax = b\).
    \item There exists \(y \in \bF^m\) such that \(yA \geq 0\) and \(y b < 0\).
  \end{itemize}
\end{lemma}

That is, either the system \(Ax = b\) can be solved, or this system is incoherent in the sense that aggregating the equalities with weights \(y\) gives a contradiction.

\mnote{
One can show this version of Farkas' Lemma by applying the geometric version to \(\tilde{A} = [A, -A]\) and \(\tilde{b} = b\).
}




\section{The No-Trade Theorem}
Voluntary trading between individual is possible only when the agents are different in some way.
One agent, for instance, may have more apples than the other (and thus a lower marginal utility for consuming apples), or they might like apples more.
Can asymmetric information be the cause of a trade?
% \Textcite{milgrom1982InformationTradea} answered this question in the negative.
We will answer this question in the negative following \Textcite{morris1994TradeHeterogeneous}.\mnote{\Textcite{morris1994TradeHeterogeneous} proved this result in a more general setting. What we present here is a special case referred to as the ``no betting case.''
I first learned this result from Ben Brooks, without whom I would not have stumbled upon this proof in the appendix of this little known paper.
}

The setup of our model is as follows:
\begin{itemize}
  \item A finite set of states \(\Theta\).
  \item A type space \(T = \prod_i T_i\) modelling information.
    % Each agent \(i\) has some prior belief \(\pi_i\) over the states and types.
    % After observing her type \(t_i\), agent \(i\) updates her prior via Bayes' rule to \(\pi_i | t_i\).
  \item Agent \(i = 1, \dots, n\).
    Agent \(i\) has prior \(\pi_i \in \Delta(T \times \Theta)\) and forms posterior \(\pi_i(t_{-i}, \theta | t_i)\) upon observing her type \(t_i\).
  \item Assume each agent is risk neutral and has quasilinear utility.
\end{itemize}

A trade is a function \(\gamma: T \times \Theta \to \bR^n\) specifying transfers to each type given the state.
Say a trade \(\gamma\) is \vocab{feasible} if no outside resource is need to facilitate it:
\[
  \sum_i \gamma_i(t, \theta) \leq 0, \quad \forall (t, \theta).
\]
Say a trade is \vocab{acceptable} if each type has a nonnegative interim payoff
\[
  \sum_{t_{-i}, \theta} \pi_i(t_{-i}, \theta | t_i) \gamma_i(t, \theta) \geq 0, \quad \forall i, t_i,
\]
with strict inequality for some \(i, t_i\).

\begin{remark}
  The acceptability condition really confused me when I initially studied this model.
  If each agent only has one \emph{realized} type, why do we care that they would also accept this trade at other types?

  One answer is that capturing the outcomes for all possible types (even the ones that can never be realized) is convenient for analysis.
  To see this, suppose there are two types of player \(A\), \(t_0^A\) and \(t_1^A\), and both \(A\) and \(B\) knows that \(A\) finds \(\gamma\) acceptable iff \(A\) she is of type \(t_1^A\).
  If \(A\) is indeed of type \(t_1^A\), she would by assumption find this trade acceptable.
  However, \(B\) does not know what type \(A\) has.
  Accordingly, he evaluates the trade using his interim payoff.
  In doing so, he analyses whether a trade will be accepted at each possible type profile \(t\) and weigh the corresponding payoffs using the probabilities of the types.
  Instead of going through such trouble, we can simply impose that a trade must be acceptable at all \(t\)'s --- and doing so is without loss, since we can also set the payoffs to be \(0\) when a trade is not accepted at some profile --- so that interim payoffs has the simple form above.
\end{remark}

Given systems of inequalities described above, it is only natural to apply some form of Farkas' lemma.

\begin{lemma}\label{lem:farkas-variant-1}
  Given \(A \in \bR^{m \times n}\) and column \(j\), exactly one of the following is true:
  \begin{enumerate}[label=(\roman*)]
    \item There exists \(x \geq 0\), \(x_j > 0\), such that \(Ax = 0\).
    \item There exists \(y\) such that \(yA \geq 0\) with strict inequality for column \(j\).
  \end{enumerate}
\end{lemma}
\begin{proof}
  Apply \Cref{lem:farkas-geom} to \(\tilde{A} = [A', e_j']'\) and \(\tilde{b} = e_{m + 1}\).
\end{proof}

\begin{lemma}\label{lem:farkas-variant-2}
  Given \(A \in \bR^{m \times n}\) and a set of columns \(J\), exactly one of the following is true:
  \begin{enumerate}[label=(\roman*)]
    \item There exists \(x \geq 0\), \(x_j > 0\) for all \(j \in J\), such that \(Ax = 0\).
    \item There exists \(j \in J\) and \(y \in \bR^m\) such that \(yA \geq 0\) for all \(j\), with strict inequality for each column \(j\).
  \end{enumerate}
\end{lemma}
\begin{proof}
  From \Cref{lem:farkas-variant-1}, (ii) implies not (i).
  On the other hand, if (ii) does not hold, then \Cref{lem:farkas-variant-1} implies that for each \(j \in J\), there exists \(x^j \geq 0\) with \(x_j^j > 0\), such that \(Ax^j = 0\).
  Then \(x \coloneq \sum_j x^j\) satisfy (i).
\end{proof}

We can now state and prove a version of the no-trade theorem:
\begin{theorem}
  There is no trade if and only if \(\pi_i\) are the same across \(i\).
\end{theorem}
\begin{proof}
  A trade exists if and only if there are \(\gamma_i\) for each \(i\) such that
  \begin{align*}
    - \sum_i \gamma_i(t, \theta)
    &\geq 0, \quad \forall t, \theta \\
    \sum_{t_{-i}, \theta} \pi_i(t_{-i}, \theta | t_i) \gamma_i(t, \theta)
    &\geq 0, \quad \forall i, t_i,
  \end{align*}
  with the second line being strict for some \(t_i\).
  This maps into \Cref{lem:farkas-variant-2} case (ii), with rows indexed by \(T \times \Theta \times \{1, \dots, n\}\), columns indexed by \((T \times \Theta) \sqcup (\sqcup_i T_i)\), and \(J = \sqcup_i T_i\).
  Thus there is no trade if and only if there exists \(\tilde{\phi}: T \times \Theta \to \bR_+\) and \(\tilde{\lambda}_i : T_i \to \bR_{++}\) such that
  \[
    - \tilde{\phi}(t, \theta)
    + \tilde{\lambda}_i(t_i) \pi_i(t_{-i}, \theta | t_i)
    = 0.
  \]
  In this case, \(\tilde{\phi} / \sum \tilde{\phi}\) is a common prior.
\end{proof}




\begin{itodo}
  Add extension to general setting?
  Talk a sentence or two about milgrom-stokey.  
\end{itodo}

